\section{Aquinas on Trade with Foreigners}

\begin{quotex}
For the country which needs considerable imports for its support must tolerate continuous intercourse with foreigners ... who, having been brought up under different laws and customs, behave in many way differently from the inhabitants of the country, so that these latter are spurred on to act similarly, and social life is disturbed. Again, if the citizens devote their lives to trade, the way will be opened to many vices. For, since the aim of traders is especially to make money, familiarity with trade leads to the awakening of greed in the hearts of the citizens. The result is that everything in the State will be put up for sale, mutual confidence will be destroyed and an atmosphere favourable to deceit and fraud created. Everyone, growing careless about the Common Good, will seek only his own advantage. The cultivation of virtue will decline, since honour, the reward of virtue, will be bestowed indiscriminately upon all comers. Hence, in such a State, social morality will inevitably suffer.
\flright{\textsc{St. Thomas Aquinas}, \emph{The Governance of Princes}}
\end{quotex}

It is clear that Aquinas is criticizing making trade and economics the primary end of society. Even the bishops today, ignoring their Thomist heritage, see everything in economic terms and ignore the resulting social disruption. For Aquinas, on the other hand, the social disruption arising from the ``continuous intercourse with foreigners'' is the most important consideration. This shows how far the Vatican II Church has moved from Tradition ... apart from their peculiar animus against abortion, they reveal themselves to be \emph{modern} men in every other regard.

It is tempting to say that the quoted passage is pure prophecy, as it has all come to pass. However, it is really the result of clear and rational thought.


\flrightit{Posted on 2008-08-28 by Cologero}

\begin{center}* * *\end{center}

\begin{footnotesize}
\begin{sffamily}
\texttt{Arturo Vasqiez on 2010-05-21 at 20:19 said:}

Then America must be screwed, since we are all immigrants here.

\hfill

\texttt{Cologero on 2010-05-22 at 11:37 said:}

Mr. Vasqiez, this blog takes no position on political issues, unless at some future date, they align with the proper direction of the world process. Also, as philosophers, our individual personal situations are of no concern; circumstances have brought us where we are, so that is where we must begin. In the Kali Yuga, our task is to deal with our situations and circumstances, however deviant from the idea.

Nevertheless, that does not absolve us from the task to understand our situations and circumstances dispasstionately and intellectually. So a careful reading of Aquinas shows he puts not blame on ``immigrants'', but rather on the greed for money and the subsequent social disruptions that arise from it.

In the case of America, these disruptions affect the immigrants as much as the original populations, since the former are both economially exploited and deprived of their social networks. An unfortunate side-effect is the re-emergence of tribalism.

The economic and tribal factors mitigate against the notion of the common good, which is the moral principle for a nation.

The economic motive prefers the particular good over the common good.

The tribal motive preferes the good of the subgroup over the common good.

This is undeniable and if anything will ``screw'' America, that will be the cause, not immigrants per se.

Although we see the clan, tribe, nation as natural human developments, we do not regard them as absolute. Rather we are in favour of Empire, or cosmopolitanism, under the guiding spirit, with nations and tribes hierarchically arranged. This will have to be left as a topic for a future post, but it will be in the direction of the Roman Empire under the man-god Caesar, and the Holy Roman Empire under the God-man.


\hfill

\end{sffamily}
\end{footnotesize}

